\section{综合讨论与研究局限}

\subsection{表示设计与融合增量}

对 RQ1，主要的误差下降出现在从基础 ExtraTrees 改为丰富特征树的阶段，说明在此任务中优先改进树分支具有价值。三折筛查在固定树参数和分组的条件下提供了更具体的证据：加入显式交互和采用母体加权时，观察到的误差均向下变化。五折评价中的丰富树会逐折选择特征和加权配置，因此整体改善包括这些选择的共同作用。

对 RQ2，TabPFN 进一步降低了平均误差，但融合收益小于从基础树到丰富树的变化。不同外层折选出的权重及母体间的误差差异说明，补充分支的价值取决于样本组成。复杂度较低的ExtraTrees 单模型适合作为快速回归基准，当前双分支适合需要同时查看两种学习器预测的条件筛选。模型流程经过两轮开发比较后确定，锁定方案后在独立来源数据上验证，才能评价这一增量能否重复。

\subsection{预测依据与推广范围}

对 RQ3，特征归因和温度扫描说明了部署模型如何利用结构与条件信息。组合内实测分析提供另一种证据，即同一完整药物--溶剂配对内部，温度与测量值如何共同变化。两类证据分别有助于检查模型行为和理解数据构成。记录缺少测量规程和样品状态等元数据，温度的干预效应需要有可比条件的测量才能估计。

母体分组验证直接对应已观测数据来源中的未见母体预测。推广到未见骨架时可采用骨架分组，推广到未见溶剂时需单独留出溶剂，超出观测温区的可靠性应通过该温区的实测记录核验。低溶解度记录的系统性高估会影响筛选排序，宜在独立评价中优先检查这一目标区间，并用独立于最终评价的样本拟合和验证校准方案。
